\subsection{Retrieval and Evaluation Software}
\label{app:implementation-software}
\paragraph{Retrieval and reranking.}
The SRA BM25 implementation~\citep{robertson2009bm25} uses NumPy/SciPy sparse operations with $k_1=1.5$ and $b=0.75$. Its tokenizer lowercases text, replaces non-word/non-whitespace characters with spaces, and splits on whitespace. Vector retrieval uses Qwen3-Embedding-0.6B~\citep{zhang2025qwen3embedding} with L2-normalized embeddings and cosine scoring. Its native encoder uses batch size one and a 32,768-token document limit; the query is an instruction followed by the verbatim task question. SkillSight uses the same embedding-model family with its background-calibrated dense and lexical scoring implementation~\citep{xiao2026skillsightcalibratinggenericcontent}. Its adapter uses a 300-candidate internal pool and scoring batches of 64 before returning the requested top-$k$ ranking. The routing interface preserves this ranking with scores $1/(1+r)$ for zero-based rank $r$, rather than interpreting the fused scores as calibrated probabilities.

SkillRouter uses the released 0.6B reranker~\citep{zheng2026skillrouter}. Its input retains the query and joins each skill's name, first 500 description characters, and first 2,000 content characters. The implementation uses batch size one, a 4,096-token total input limit, left padding, evaluation mode without gradients, and bfloat16 on CUDA or float32 on CPU. Candidates are ordered by the difference between the \texttt{yes} and \texttt{no} logits, without fusion with the original retrieval score.

\paragraph{Skill-Pro.}
Skill-Pro represents each skill using an activation condition, an execution procedure, and a termination condition~\citep{pmlr-v306-mi26d}. In our implementation, we freeze the existing skills and generate only a natural-language activation condition for each skill. This condition describes when the corresponding skill should be invoked; the skill content remains unchanged.

\paragraph{SkillDAG.}
We evaluate SkillDAG with its default settings~\citep{bai2026skilldag}. Starting from the five highest-ranked retrieved skills, we expand the skill graph and rank the resulting skills by graph depth from these seeds and retrieval score. We retain at most the top-$k$ ranked skills, truncating the selection if expansion produces more than $k$ candidates. For downstream evaluation, we use $k=10$.

\paragraph{Task evaluation.}
We use benchmark-specific SRA task evaluators. LogicBench compares normalized extracted labels. ToolQA uses normalized exact match, with a numerical fallback requiring absolute error below $10^{-6}$ and a Boolean-label fallback. MedCalc-Bench applies output-type-specific rules: date and gestational-age equality, rounded integer comparison, or inclusive supplied bounds for decimal answers. CHAMP uses canonical, numerical, and SymPy-based equivalence checks; TheoremQA uses type-aware numerical and list comparisons with \texttt{latex2sympy2} parsing. Their shared numerical-tolerance helper uses 4\% relative tolerance with a $10^{-9}$ absolute floor. BigCodeBench extracts code with tree-sitter and executes the supplied unit tests, with address-space and data limits of 30\,GiB and a stack limit of 10\,MiB.

\paragraph{Examples and split construction.}
ToolQA's ReAct prompt~\citep{yao2023react} contains eight fixed easy and three fixed hard demonstrations, in that order. These are separate from the training-history examples used for predicate repair in Appendix~\ref{app:procedure-replay}. The downstream runner defaults to this ReAct path for ToolQA and direct-answer inference for the other tasks. The skill-balanced split was constructed with SciPy 1.18.1 and its HiGHS mixed-integer solver, using normalized absolute-deviation objectives and per-skill split-coverage constraints as described in Appendix~\ref{app:data-splits}.
